% Local statements restated from work/degree6/dyadic_hecke_kernel.tex,
% Lemma lem:dyadic-Hecke, and work/degree6/quotient_cartier_criterion.tex,
% Lemma lem:quotient-Cartier. The target below is the full Jacobian.
\begin{lemma}\label{lem:decic49-dyadic-kernel}
Let $A/\mathbb Z_2$ be an abelian scheme, and let $T,\delta$ be
endomorphisms satisfying $T_{\mathbb F_2}=F+\delta_{\mathbb F_2}V$.
Then $T-\delta$ kills the torsion kernel of reduction in $A(\mathbb Q_2)$.
\end{lemma}
\begin{proof}
The torsion kernel has exponent two. Indeed, multiplication by two
in the formal group sends parameters in $2\mathbb Z_2$ into
$4\mathbb Z_2$, and the latter subgroup is torsion free by the
formal logarithm.
Let $z$ be a nonzero point in this kernel. The schematic closure
of $\{0,z\}$ in $A[2]$ is finite flat of rank two with connected
special fibre. It is $\mu_2$: in the rank-two Tate--Oort description
\cite{TO}, the parameters satisfy $ac=2$; connected special fibre
forces $a\in2\mathbb Z_2$, so $c$ is a unit, giving $\mu_2$ after
rescaling. Write $\iota:\mu_2\hookrightarrow A[2]$ for its inclusion.
Naturality and $F_{\mu_2}=0$, $V_{\mu_2}=1$ show that
$(T-\delta)\overline\iota=0$ on the special fibre.
Reduction on $\operatorname{Hom}_{\mathbb Z_2}(\mu_2,A[2])$ is injective.
By Cartier duality, this is the corresponding assertion for maps
from $A^\vee[2]$ to the constant etale group $\mathbb Z/2\mathbb Z$.
Their open and closed fibres are determined by their special
fibres over the henselian ring $\mathbb Z_2$.
Thus $(T-\delta)\iota=0$, proving the assertion.
\end{proof}

\begin{lemma}\label{lem:decic49-Cartier}
Let $Y/\mathbb Q$ be a modular curve of odd level with good
reduction at $2$, rational cusp $c_0$, and Jacobian $J$ of rank zero
over $\mathbb Q$. Suppose that its cusp scheme is etale at $2$
and that $T_2=F+\delta V$ on the special fibre, where
$\delta=\langle2\rangle$. Put
$U=(\mathcal C-1)H^0(Y_{\mathbb F_2},\Omega^1)$.
If restriction from $U$ to an effective cuspidal divisor
$\overline C$ of degree $d$ is surjective, then no noncuspidal
rational effective divisor of degree $d$ reduces to $\overline C$.
\end{lemma}
\begin{proof}
Let $C$ be the unique cuspidal lift of $\overline C$ over
$\mathbb Z_2$. This is a $\mathbb Q_2$-rational divisor: a closed
cusp is lifted together with its entire Frobenius orbit.
Suppose that $D$ is a noncuspidal rational effective divisor
reducing to $\overline C$. The classes $[D-dc_0]$ and $[C-dc_0]$
are torsion, respectively by rank zero over $\mathbb Q$ and
Manin--Drinfeld. Consequently $[D-C]\in J(\mathbb Q_2)$ is
torsion and reduces to zero. Lemma~\ref{lem:decic49-dyadic-kernel}
gives $(T_2-\delta)[D-C]=0$.

Use the full Jacobian as target in
\[
 f:Y^{(d)}\longrightarrow J,\qquad
 E\longmapsto(T_2-\delta)[E-dc_0].
\]
The cotangent space of $J_{\mathbb F_2}$ is canonically
$H^0(Y_{\mathbb F_2},\Omega^1)$. On it, Frobenius pulls back to
zero and Verschiebung pulls back to Cartier. Hence
\[
 (T_2-\delta)^*=\delta^*(\mathcal C-1),
\]
whose image is $U$, since diamonds commute with Cartier.
The Abel--Jacobi cotangent map at $\overline C$ is restriction to
$H^0(\overline C,\Omega^1|_{\overline C})$, including when
$\overline C$ is nonreduced. Thus the stated rank condition makes
$f$ a formal immersion there. This uses the full Jacobian and its
canonical cotangent identification; no cotangent exact sequence
for a quotient in characteristic two is needed.
Both $D$ and $C$ give $\mathbb Z_2$-sections, with the same
reduction and the same image under $f$. Formal immersion implies
$D=C$, a contradiction.
\end{proof}
